\section{Conclusion}
\label{sec:conclusion}

We evaluated explainability-guided electrode reduction for subject-independent sEMG gesture recognition using cross-fitted attribution and decision rules fixed in advance. On 43 GRABMyo participants, a 12-of-16 electrode subset was non-inferior to the full array within a 0.02 Macro-F1 engineering tolerance; smaller subsets were not. Cross-fitted TreeSHAP rankings were stable across people. They beat random subsets and were faithful under deletion with retraining. A stability penalty added nothing, and a data-free, evenly spread layout performed about as well as the attribution ranking at most sizes. The results support moderate, not aggressive, electrode reduction for population-level devices. In an exploratory re-analysis with bipolar pairs as the unit, no pair could be removed within the same tolerance. They also suggest that the benefit of attribution lies in confirming which positions matter, rather than in finding sparse layouts that simple coverage rules would miss.
